\questao{128}{corC}{Ondas e óptica}{Espalhamento da luz na neblina}{Só de conceito}
\begin{enunciado}
% gerado por ferramentas/md2tex_enunciado.py a partir de questoes/q128.md (não editar à mão)
Ao trafegar à noite, é necessário manter os faróis dos veículos acesos. Quando a umidade relativa do ar é baixa, o motorista consegue ter boa visibilidade da pista. No entanto, em noites em que há neblina, essa visibilidade é reduzida.\par
Nessa situação, a visibilidade reduzida deve-se à\par
\begin{alternativas}
  \item dispersão da luz pelas gotículas de água suspensas no ar.
  \item polarização da luz pelas gotículas de água dispersas no ar.
  \item condensação das gotículas de água na superfície do farol.
  \item solubilização das gotículas de água presentes no ar.
  \item redução da potência da luz emitida pelo farol.
\end{alternativas}

\end{enunciado}

\quer{por que o motorista enxerga pior a pista quando há neblina.}

\begin{ideia}
Para você enxergar a pista à noite, a luz do farol precisa fazer um caminho: sair do farol, chegar até a pista (ou a um
obstáculo), bater e \textbf{voltar para os seus olhos}.

A neblina é formada por \textbf{gotículas de água muito pequenas, suspensas no ar}. Quando a luz do farol encontra essas
gotículas, cada uma desvia um pouco de luz \textbf{para todos os lados}. É isso que a alternativa chama de
\textbf{dispersão} (ou espalhamento) da luz.

Duas coisas ruins acontecem ao mesmo tempo:
\begin{itemize}
  \item \textbf{menos luz chega à pista}, porque boa parte se perde pelo caminho, desviada para os lados;
  \item \textbf{parte da luz volta para o motorista}, formando aquela ``parede branca'' na frente do carro.
\end{itemize}

\textbf{Você já viu isso:} o feixe de uma lanterna ou do projetor do cinema fica visível quando há fumaça ou poeira no ar.
Só dá para ver o feixe de lado porque as partículas espalham a luz na direção dos seus olhos.

\begin{center}
\begin{tikzpicture}[font=\small]
  \node[draw=tinta,fill=cinzaclaro,rounded corners=1mm,minimum width=13mm,minimum height=7mm] (f) at (0,0) {farol};
  \draw[seta,corB] (f.east) -- (3.0,0);
  \foreach \x/\y in {3.2/0, 4.2/0.38, 4.9/-0.32, 5.6/0.2}{\fill[corA!50] (\x,\y) circle (2.4pt);}
  \foreach \a in {40,100,160,220,290,340}{\draw[corB!80,line width=0.6pt,-{Latex[length=1.4mm]}] ($(3.2,0)+(\a:0.16)$) -- ($(3.2,0)+(\a:0.8)$);}
  \draw[corB!45,line width=0.9pt,-{Latex[length=1.6mm]},dashed] (6.15,0) -- (8.6,0);
  \node[draw=tinta,fill=cinzaclaro,rounded corners=1mm,minimum width=13mm,minimum height=7mm] at (9.5,0) {pista};
  \node[cinza,anchor=north] at (4.4,-0.95) {gotículas de água: espalham a luz para todos os lados};
  \node[cinza,anchor=south] at (7.4,0.12) {chega pouca luz};
\end{tikzpicture}
\end{center}
\end{ideia}

\passoapasso
\npasso{Descubra o que muda entre a noite seca e a noite com neblina.}
O farol é o mesmo. O motorista é o mesmo. A única diferença está \textbf{no ar}: com neblina, há gotículas de água entre o
carro e a pista.

\npasso{Pense no que as gotículas fazem com a luz.}
Elas não apagam a luz nem mudam o farol. Elas \textbf{desviam} a luz do caminho reto, espalhando-a em todas as direções.

\npasso{Escolha a alternativa que descreve isso.}
É a alternativa A: dispersão da luz pelas gotículas de água suspensas no ar.

\resposta{A}{As gotículas da neblina espalham a luz do farol em todas as direções. Chega menos luz à pista e parte dela volta para os olhos do motorista.}

\begin{outras}
\alt{B} Polarizar a luz é ``filtrar'' a direção em que ela vibra (como fazem alguns óculos de sol). Não é isso que a neblina faz, e isso não explicaria a parede branca.\par
\alt{C} O problema não está no vidro do farol, e sim no ar entre o carro e a pista.\par
\alt{D} ``Solubilizar'' é dissolver uma substância em outra. As gotículas continuam lá, como gotículas: nada se dissolveu.\par
\alt{E} A lâmpada do farol emite a mesma potência com ou sem neblina. O que muda é o caminho da luz depois que ela sai.
\end{outras}

\begin{pegadinha}[Onde muita gente erra]
\begin{itemize}
  \item \textbf{Marcar C porque neblina ``lembra'' condensação.} A alternativa fala de gotículas \emph{na superfície do farol}.
        O problema está no ar, entre o carro e a pista.
  \item \textbf{Descartar A por achar que dispersão é só o arco-íris.} Em Óptica, ``dispersão'' também é o nome da separação da
        luz branca em cores. Nesta questão, a palavra foi usada no sentido de \textbf{espalhamento}.
\end{itemize}
\end{pegadinha}

\begin{guardar}[Para guardar (e usar na estrada)]
Na neblina, o \textbf{farol alto piora} a visão: ele joga mais luz nas gotículas à altura dos olhos, e mais luz volta para o
motorista. Por isso se usa farol baixo ou farol de neblina, que iluminam perto do chão.
\end{guardar}
